\documentclass{article}
\usepackage{amsmath,amssymb}
\usepackage{xurl}
\usepackage[hidelinks]{hyperref}
\setlength{\emergencystretch}{2em}
% Shared commands for the notes in docs/paper-gaps/.
% Every note loads this file via % Shared commands for the notes in docs/paper-gaps/.
% Every note loads this file via % Shared commands for the notes in docs/paper-gaps/.
% Every note loads this file via \input{command}. Project-wide notation
% belongs here, never in per-note redefinitions.

% --- Project configuration ------------------------------------------------
\newcommand{\projectIssueBase}{https://github.com/LionSR/QICLean/issues/}
\newcommand{\projectPrBase}{https://github.com/LionSR/QICLean/pull/}
\newcommand{\projectDocBase}{https://sirui-lu.com/QICLean/docs/}
% ---------------------------------------------------------------------------

\DeclareUnicodeCharacter{2080}{\ifmmode{}_{0}\else\textsubscript{0}\fi}
\DeclareUnicodeCharacter{2081}{\ifmmode{}_{1}\else\textsubscript{1}\fi}
\DeclareUnicodeCharacter{2082}{\ifmmode{}_{2}\else\textsubscript{2}\fi}
\DeclareUnicodeCharacter{2099}{\ifmmode{}_{n}\else\textsubscript{n}\fi}

\newcommand{\Lean}{\textsc{Lean}}

% A formal declaration name, upright sans, breakable at underscores.
\ExplSyntaxOn
\NewDocumentCommand{\leanid}{m}{
  \ifmmode
    \textsf{\detokenize{#1}}
  \else
    \group_begin:
    \urlstyle{sf}
    \tl_set:Nn \l_tmpa_tl {#1}
    \tl_replace_all:Nnn \l_tmpa_tl {\_} {_}
    \exp_args:NV \path \l_tmpa_tl
    \group_end:
  \fi
}
\ExplSyntaxOff

% Traceability links; use in footnotes, not in the running prose.
\newcommand{\ghissue}[1]{\href{\projectIssueBase#1}{GitHub issue \##1}}
\newcommand{\ghpr}[1]{\href{\projectPrBase#1}{PR \##1}}
\newcommand{\doclink}[1]{\href{\projectDocBase#1.html}{\path{#1}}}

% Domain notation, using \providecommand so a note-local refinement stays
% possible.
\providecommand{\tr}{\operatorname{tr}}

% Dirac notation and the identity operator, as in blueprint/src/macros.
\usepackage{dsfont}
\providecommand{\ket}[1]{|#1\rangle}
\providecommand{\bra}[1]{\langle #1|}
\providecommand{\braket}[2]{\langle #1 | #2 \rangle}
\providecommand{\ketbra}[2]{|#1\rangle\!\langle #2|}
\providecommand{\Id}{\mathds{1}}
\providecommand{\one}{\mathds{1}}
\providecommand{\C}{\mathbb{C}}
\providecommand{\MN}[1]{M_{#1}(\C)}
\providecommand{\MD}{M_D(\C)}

% Verdict marker: \gapnote{<kind>}{<status>}. Kind is the policy
% classification (clarification, local-correction, scope-restriction,
% unfaithful, false-source, open-gap); status is open, wip (elimination
% underway), resolved, or historical. Machine-read by the site index;
% prints a verdict line.
\newcommand{\gapnote}[2]{\par\noindent\textbf{Verdict:} #1 (#2).\par}
. Project-wide notation
% belongs here, never in per-note redefinitions.

% --- Project configuration ------------------------------------------------
\newcommand{\projectIssueBase}{https://github.com/LionSR/QICLean/issues/}
\newcommand{\projectPrBase}{https://github.com/LionSR/QICLean/pull/}
\newcommand{\projectDocBase}{https://sirui-lu.com/QICLean/docs/}
% ---------------------------------------------------------------------------

\DeclareUnicodeCharacter{2080}{\ifmmode{}_{0}\else\textsubscript{0}\fi}
\DeclareUnicodeCharacter{2081}{\ifmmode{}_{1}\else\textsubscript{1}\fi}
\DeclareUnicodeCharacter{2082}{\ifmmode{}_{2}\else\textsubscript{2}\fi}
\DeclareUnicodeCharacter{2099}{\ifmmode{}_{n}\else\textsubscript{n}\fi}

\newcommand{\Lean}{\textsc{Lean}}

% A formal declaration name, upright sans, breakable at underscores.
\ExplSyntaxOn
\NewDocumentCommand{\leanid}{m}{
  \ifmmode
    \textsf{\detokenize{#1}}
  \else
    \group_begin:
    \urlstyle{sf}
    \tl_set:Nn \l_tmpa_tl {#1}
    \tl_replace_all:Nnn \l_tmpa_tl {\_} {_}
    \exp_args:NV \path \l_tmpa_tl
    \group_end:
  \fi
}
\ExplSyntaxOff

% Traceability links; use in footnotes, not in the running prose.
\newcommand{\ghissue}[1]{\href{\projectIssueBase#1}{GitHub issue \##1}}
\newcommand{\ghpr}[1]{\href{\projectPrBase#1}{PR \##1}}
\newcommand{\doclink}[1]{\href{\projectDocBase#1.html}{\path{#1}}}

% Domain notation, using \providecommand so a note-local refinement stays
% possible.
\providecommand{\tr}{\operatorname{tr}}

% Dirac notation and the identity operator, as in blueprint/src/macros.
\usepackage{dsfont}
\providecommand{\ket}[1]{|#1\rangle}
\providecommand{\bra}[1]{\langle #1|}
\providecommand{\braket}[2]{\langle #1 | #2 \rangle}
\providecommand{\ketbra}[2]{|#1\rangle\!\langle #2|}
\providecommand{\Id}{\mathds{1}}
\providecommand{\one}{\mathds{1}}
\providecommand{\C}{\mathbb{C}}
\providecommand{\MN}[1]{M_{#1}(\C)}
\providecommand{\MD}{M_D(\C)}

% Verdict marker: \gapnote{<kind>}{<status>}. Kind is the policy
% classification (clarification, local-correction, scope-restriction,
% unfaithful, false-source, open-gap); status is open, wip (elimination
% underway), resolved, or historical. Machine-read by the site index;
% prints a verdict line.
\newcommand{\gapnote}[2]{\par\noindent\textbf{Verdict:} #1 (#2).\par}
. Project-wide notation
% belongs here, never in per-note redefinitions.

% --- Project configuration ------------------------------------------------
\newcommand{\projectIssueBase}{https://github.com/LionSR/QICLean/issues/}
\newcommand{\projectPrBase}{https://github.com/LionSR/QICLean/pull/}
\newcommand{\projectDocBase}{https://sirui-lu.com/QICLean/docs/}
% ---------------------------------------------------------------------------

\DeclareUnicodeCharacter{2080}{\ifmmode{}_{0}\else\textsubscript{0}\fi}
\DeclareUnicodeCharacter{2081}{\ifmmode{}_{1}\else\textsubscript{1}\fi}
\DeclareUnicodeCharacter{2082}{\ifmmode{}_{2}\else\textsubscript{2}\fi}
\DeclareUnicodeCharacter{2099}{\ifmmode{}_{n}\else\textsubscript{n}\fi}

\newcommand{\Lean}{\textsc{Lean}}

% A formal declaration name, upright sans, breakable at underscores.
\ExplSyntaxOn
\NewDocumentCommand{\leanid}{m}{
  \ifmmode
    \textsf{\detokenize{#1}}
  \else
    \group_begin:
    \urlstyle{sf}
    \tl_set:Nn \l_tmpa_tl {#1}
    \tl_replace_all:Nnn \l_tmpa_tl {\_} {_}
    \exp_args:NV \path \l_tmpa_tl
    \group_end:
  \fi
}
\ExplSyntaxOff

% Traceability links; use in footnotes, not in the running prose.
\newcommand{\ghissue}[1]{\href{\projectIssueBase#1}{GitHub issue \##1}}
\newcommand{\ghpr}[1]{\href{\projectPrBase#1}{PR \##1}}
\newcommand{\doclink}[1]{\href{\projectDocBase#1.html}{\path{#1}}}

% Domain notation, using \providecommand so a note-local refinement stays
% possible.
\providecommand{\tr}{\operatorname{tr}}

% Dirac notation and the identity operator, as in blueprint/src/macros.
\usepackage{dsfont}
\providecommand{\ket}[1]{|#1\rangle}
\providecommand{\bra}[1]{\langle #1|}
\providecommand{\braket}[2]{\langle #1 | #2 \rangle}
\providecommand{\ketbra}[2]{|#1\rangle\!\langle #2|}
\providecommand{\Id}{\mathds{1}}
\providecommand{\one}{\mathds{1}}
\providecommand{\C}{\mathbb{C}}
\providecommand{\MN}[1]{M_{#1}(\C)}
\providecommand{\MD}{M_D(\C)}

% Verdict marker: \gapnote{<kind>}{<status>}. Kind is the policy
% classification (clarification, local-correction, scope-restriction,
% unfaithful, false-source, open-gap); status is open, wip (elimination
% underway), resolved, or historical. Machine-read by the site index;
% prints a verdict line.
\newcommand{\gapnote}[2]{\par\noindent\textbf{Verdict:} #1 (#2).\par}

\title{The constants of the bootstrapping step for degenerate ground spaces}
\author{}
\date{2026-10-26}
\begin{document}
\maketitle
\gapnote{false-source}{open}

The bootstrapping estimate of Arad, Landau, Vazirani and Vidick is stated in
Proposition \texttt{prop:albootstrap-dg} of
\href{https://arxiv.org/abs/1602.08828}{arXiv:1602.08828v2}, source lines
1062--1075. It bounds the minimal dimension $q$ of a $0.015$-viable subspace for
the ground space $T$. Its proof takes $W$ of minimal dimension, observes that
$W$ is $0.04$-viable for the subspace $\tilde T$ of Theorem \texttt{t:ais},
samples a random $W'\subseteq W$ of dimension $s=q/(2D^2)$, applies the
operators $A_1,\dots,A_{D^2}$, and contradicts minimality. Within the interval
setting, three quantitative details need adjustment; none changes the dimension
bound. The region restriction is explained below.

\emph{The target of the sample.} Line 1068 applies Lemma \texttt{l:rs} to
obtain a sample viable for $T$, but the next step applies the $A_i$, which
improve viability for $\tilde T$. The sample must be viable for $\tilde T$,
which holds because $W$ is $0.04$-viable for $\tilde T$.

\emph{The sampling condition.} The condition
\texttt{eq:bootstrap-cond-dg} (line 1070) is stated as
$s=\Omega(\log q+r\log(q/s))$, which is meant to make the failure probability
of Lemma \texttt{l:rs} smaller than one. The formal statement uses that
requirement directly, with the failure probability
$\eta(q)=(1+2/\nu)^{2r}q\,e^{-s/16}$, $\nu=\sqrt{0.96\,s/(8q)}$, of the
corrected Lemma \texttt{l:rs} (note \path{alvv16_random_sampling_net_count}),
and with $s=\lfloor q/(2D^2)\rfloor$.

\emph{The constants.} Line 1074 states $(32D^2)^2\Delta+0.005\le0.1\le0.15$
and concludes a contradiction with the threshold $0.015$. Robustness,
Lemma \texttt{l:close}, doubles the error, so the span of the $A_iW'$ is
$2(\Delta/\delta'^2+0.005)$-viable for $T$, where $\delta'=0.96\,s/(8q)$ is the
overlap of the sample. The contradiction therefore needs
$\Delta/\delta'^2\le0.0025$. For $q\ge4D^2$ we have $s\ge q/(4D^2)$ and
$\delta'\ge0.03/D^2$, so $\Delta/\delta'^2\le D^4\Delta/0.0009$, and
$D^{12}\Delta\le10^{-5}$ gives $D^4\Delta\le10^{-5}/D^8$. This is at most
$0.0025\cdot0.0009$ for $D\ge2$, but not for $D=1$. Theorem \texttt{t:ais}
has $D=e^{\tilde O(\gamma^{-1}\log^3d)}\ge2$, so the formal statement assumes
$D\ge2$ and $q\ge4D^2$.

The formal statements are
\leanid{Submodule.exists_isViable_finrank_le_half} for one step and
\leanid{Submodule.exists_isViable_finrank_lt_of_bootstrap} for the
proposition, which replaces the minimality argument by repeated halving
from the whole factor. The explicit threshold $rD^{24}$, with
$r=\dim T'\ge1$, is supplied by
\leanid{Submodule.exists_isViable_finrank_lt_rank_mul_pow} under the same
improving-family assumptions. The physical interval construction remains
a separate application.

\section{The region in the physical proposition}

Proposition \texttt{prop:albootstrap-dg}, source line 1063, permits an
arbitrary subset $J\subseteq\{1,\ldots,n\}$. Its proof at line 1067 invokes
Theorem \texttt{t:ais} with middle region $M=J$, although that theorem,
line 1010, assumes three contiguous blocks. The asserted dimension bound
cannot hold for arbitrary subsets.

Indeed, take $n=2m$ qubits and put one Bell-pair penalty on each disjoint
adjacent pair, with all intervening interactions zero:
\begin{align}
    H&=\sum_{j=1}^m\bigl(I-|\Phi^+\rangle\langle\Phi^+|\bigr)_{2j-1,2j},
    &|\Phi^+\rangle&=\frac{|00\rangle+|11\rangle}{\sqrt2}.
    \label{eq:bootstrap_disconnected_bell_hamiltonian}
\end{align}
Every local term is a positive contraction. The terms act on disjoint
pairs, so their common kernel is the unique ground state
$|\psi\rangle=|\Phi^+\rangle^{\otimes m}$, and the gap is one.
Thus Assumption (DG), source lines 267 and 275--279, holds with
$d=2$, $r=1$, $\varepsilon_0=0$ and $\gamma=1$.
For $J=\{1,3,\ldots,2m-1\}$, the reduced density matrix is
$\rho_J=I/2^m$. Hence every subspace $W\subseteq\mathcal H_J$ satisfies
\begin{align}
    \langle\psi|(P_W\otimes I)|\psi\rangle
    &=\operatorname{tr}(\rho_JP_W)=\frac{\dim W}{2^m}.
    \label{eq:bootstrap_disconnected_bell_overlap}
\end{align}
By the definition of viability, source lines 433--439, $0.015$-viability
therefore requires $\dim W\ge(197/200)2^m$. In contrast, the bound in
Proposition \texttt{prop:albootstrap-dg} is
$r\exp(\widetilde O(\gamma^{-1}\log^3d))$, as defined by the macro
\texttt{alone} at line 166. With $d=2$, $r=1$ and $\gamma=1$, this bound
is independent of $m$, a contradiction for sufficiently large $m$.

The physical proposition must explicitly restrict $J$ to an interval.
This restriction suffices for the original area-law deduction: line 1080
uses only the left and right intervals determined by a single cut.
In the example above, any single cut crosses at most one Bell pair, so
the Schmidt rank at that cut is at most two.
The restriction leaves the abstract bootstrapping results unchanged, since those
results assume an improving operator family for the specified tensor
factor. Their quantitative corrections above are resolved. A physical
interval theorem must retain its stated region restriction; it does not
formalize the printed arbitrary-subset proposition.

\section{One common event for resampling and reference trimming}

The informal merge at source lines 711--755 describes repeated
amplification followed by trimming. The formal development treats a
repeated-resampling variant of it: the sampling step of amplification is
retained literally at each round. This is a different operation from
Figure \texttt{fig:merge-prime}, lines 1664--1673, which samples once
before its loop. Figure \texttt{fig:merge-prime} and Theorem
\texttt{t:alg1} are not formalized, and nothing below identifies the two
procedures.

Let $T_0$ be the original target and $(T_x)$ a nonempty family of improving
targets. Suppose $T_0$ is $\varepsilon$-close to every $T_x$, every $T_x$ is
$1/200$-close to $T_0$, and the same family of at most $D^2$ operators
improves viability for every $T_x$. Assume $D\ge2$,
$D^{12}\Delta\le10^{-5}$, and $\varepsilon\le(0.03/D^2)^2/100$.
For a $0.04$-viable current space of dimension $q\ge4D^2$, put
$s=\lfloor q/(2D^2)\rfloor$ and $\alpha=0.96s/(8q)$.
One Borel sampling-success event for the fixed target $T_0$ has complement
measure at most $(1+2/\sqrt\alpha)^{2\dim T_0}q e^{-s/16}$.
On this event, weighted composition with parameter $\alpha/4$ makes the
same sample $(1-\alpha/2)$-viable for every $T_x$. Its single common-family
image therefore has error $4\Delta/\alpha^2$ for every $T_x$, and error
at most $27/2500=0.0108$ for $T_0$.\footnote{Formal theorem:
\leanid{Submodule.exists_measurable_half_sample_uniform_image_event}.}

A reference to $T_0$, rather than to the improving targets, suffices for
the subsequent trim. If its two directed errors are at most $10^{-5}$,
its vectors have Schmidt rank at most $b$ at the $m$ cuts, and
$\xi\sqrt{mb}\le10^{-3}$, weighted reference trimming with parameter
$1/20$ bounds the final error by
\[
\begin{aligned}
 &(21/20)^2(27/2500)+(21/20)21\cdot10^{-5}\\
 &\qquad{}+21\cdot10^{-5}+(21/10)10^{-3}
   =0.0144375<0.015.
\end{aligned}
\]
Both the common-family image and its literal collective trim have
dimension at most $q/2$ for every unitary, including unsuccessful samples.
The event is chosen for a fixed target; no assertion that an uncountable
intersection over centers is Borel is required.\footnote{Formal theorem:
\leanid{Entropy.exists_measurable_half_sample_reference_trim_event}.}

\paragraph{The terminal deterministic step.}
A stopped resampling process may take no active round. Its retained
$0.04$-viability therefore needs one final amplification and trim before
claiming $0.015$-viability. If the original target and every improving target
are mutually $1/200$-close, the current space is $9/100$-viable for every
improving target. Apply the same full family once, without another sample.
Its same image has exact error $\Delta/(91/100)^2$ for every improving
target. Under $D\ge2$ and $D^{12}\Delta\le10^{-5}$ this is at most $10^{-4}$;
no nonnegative-error premise is needed. Transfer back and the preceding
original-reference trimming budgets give final error at most
$3/200$.\footnote{Formal theorem:
\leanid{Entropy.isCloseTo_terminal_family_reference_trim}.}
Unconditionally, both the full image and its literal trim have dimension
at most $|I|\dim W\le D^2\dim W$.\footnote{Formal theorem:
\leanid{Entropy.finrank_family_image_trim_le}.}

\paragraph{The stopped resampling process.}
Under the hypotheses of the one-step event above, together with the
reference hypotheses, the stopped process takes a fresh sample whenever the
current rank is at least a threshold and otherwise keeps the current space.
There is one coefficient $K\ge1$, depending only on $D$, the exponent $a$
and $\eta>0$, and chosen before the length $N$, the target dimension bound
$r\ge\dim T_0$ and the ambient dimension $d$, with the following property.
For the threshold $K(r+1)(N+2)^2$ and any number $R$ of rounds with
$d\le2^R$, the event that every state is $0.04$-viable for $T_0$ and every
active next state is $3/200$-viable for $T_0$ is Borel for the product of
the Haar laws and has failure probability at most $R\eta/(N+2)^a$. On every
history the ranks do not increase and the final rank is below the
threshold. The initial matrix need not be an orthogonal
projection.\footnote{Formal theorem:
\leanid{Entropy.exists_canonicalHalfSampleTrimIterate_joint_probability}.}
Appending the terminal deterministic step gives a final space that is
$3/200$-viable for $T_0$ except on an event of probability at most
$R\eta/(N+2)^a$; on the same event its untrimmed image has error
$\Delta/(91/100)^2$ for every improving target, and its dimension is at most
$D^2$ times the final rank.\footnote{Formal theorem:
\leanid{Entropy.canonicalHalfSamplingParent_joint_probability_of_threshold},
for any coefficient satisfying the scalar threshold of
\leanid{Submodule.exists_bootstrapFailure_polynomial_threshold}.}
The proofs use only fixed-input Haar estimates and the dependence of each
round on the preceding seeds; no measurable selection of success events is
assumed.

\paragraph{What remains open.}
The statements above are proved for the repeated-resampling variant under
their stated family and reference hypotheses. The following are not
established by them: the sample-once procedure of Figure
\texttt{fig:merge-prime} and its analysis in Theorem \texttt{t:alg1}; the
construction, in the physical interval setting, of a common improving
family and a reference satisfying the uniform target and reference
hypotheses; and any running-time bound.

\end{document}
